% ==============================================================================
% 附：重排版校订注记（汇总原书文字误植、公式印刷瑕疵与插图数理订正）
% ==============================================================================
\clearpage
\phantomsection
\section*{重排版校订注记}
\addcontentsline{toc}{section}{重排版校订注记}
\markboth{重排版校订注记}{重排版校订注记}

{\small
\setlength{\abovedisplayskip}{3pt}
\setlength{\belowdisplayskip}{3pt}
本重排版以《形形色色的曲线》（蒋声著，上海教育出版社 1985 年第 1 版）为底本，在保留原书数学记号（如 $\tg$、$\ctg$、$\sh$、$\ch$、$\thyp$、$\arc{AB}$ 等）与总体编排风貌的前提下，对其中的排印笔误及插图数理偏差进行了逐一核查与修订。所有实质性订正除在正文相应位置以脚注随页标注外，兹汇总于下，以备读者查考。

\vspace{-1.5mm}
\subsection*{一、文字、符号与公式校订}

\begin{enumerate}\setlength{\itemsep}{0.5pt}\setlength{\parsep}{0pt}\setlength{\topsep}{2pt}
  \item \textbf{第 18 节（长幅摆线与坐标轴位置）}：原书“并且有一部分弧出现在 $y$ 轴下方”系排印笔误；据参数方程 $y = r - a\cos t$ 及图文所示，当 $a > r$ 时部分曲线位于横轴下方，已订正为“$x$ 轴下方”。
  \item \textbf{第 32 节（凸轮轮廓曲线记号）}：原书第 51 页“曲线 $C$ 上的圆弧 $\arc{ABC}$ 的圆心就是 $O$ 点”将轮廓曲线记号误排作大写 $C$；据上文“轮廓曲线 $c$”统一订正为小写 $c$。
  \item \textbf{第 58 节（麦克劳林三等分角曲线方程）}：原书第 102 页（1）式及第 103 页（2）式在排印减号时误植双短横（如 \texttt{(x--X)}），已据上下文订正为单减号 \texttt{(x-X)}。
  \item \textbf{第 78 节（曳物线图注编号）}：原书第 135 页“图 78.2”图注误印作“图 73.2”，已据节次及正文引用订正为“图 78.2”。
  \item \textbf{第 84 节（$n$ 次抛物线文字脱漏）}：原书第 151 页末句“在 $x=0$ 附近上述方程所表示的 $n$ 抛物线”漏排“次”字；据上下文及全节定义订正为“$n$ 次抛物线”。
  \item \textbf{第 93 节（对数螺线位似对应点）}：原书“因而 $OB' > OD$，从 $D$ 到 $B$ 的变换是放大”中，点 $B$ 系 $B'$ 之误植；因位似变换将 $C_1$ 上极角为 $\theta_1$ 之点 $D$ 映至 $C_2$ 上同极角之点 $B'$，已订正为“从 $D$ 到 $B'$ 的变换是放大”。
\end{enumerate}

\vspace{-1.5mm}
\subsection*{二、插图数理严谨性与绘法修订}

\begin{enumerate}\setlength{\itemsep}{0.5pt}\setlength{\parsep}{0pt}\setlength{\topsep}{2pt}
  \item \textbf{第 20 节 图 20.1（卧式旋耕机刀片轨迹）}：原书图中刀片“1”的编号位置不够清晰，且四把刀片所绘长幅摆线轨迹的起点截断参差不齐；重绘时统一自各刀尖（$t=0$）同步绘出轨迹并依序标清刀片 1～4。
  \item \textbf{第 24 节 图 24.2（$m=\frac{2}{3}$ 内摆线）}：原书图 24.2 中 $m=\frac{2}{3}$ 一图误绘为带有三个绕扣的闭合圆环。事实上，按内摆线参数方程
  \[\begin{aligned}
    x &= (R - mR)\cos mt + mR\cos(t - mt), \\
    y &= (R - mR)\sin mt - mR\sin(t - mt),
  \end{aligned}\]
  当 $m=\frac{r}{R}=\frac{2}{3}$ 时，$x'(t)=y'(t)=0$ 恰在 $t=0,\, 2\pi,\, 4\pi$ 处产生三个尖点（周期 $6\pi$），其轨迹为无绕扣的三尖内摆线（三尖瓣线，几何上与 $m=\frac{1}{3}$ 同构）。本版据正确数学方程重绘。
  \item \textbf{第 38 节 图 38.1(3)（$b=\frac{5}{3}$ 五叶玫瑰线）}：原书将五叶玫瑰线 $\rho = a\sin\frac{5}{3}\theta$ 上下倒置绘制（将竖直花瓣朝上）；按正方程当 $\theta = \frac{3\pi}{2}$ 时 $\rho = a$ 位于正下方（$270^\circ$），本版据极坐标方程真实方位重绘。
  \item \textbf{第 45 节 图 45.1（$m=\frac{1}{3}$ 刀尖轨迹）}：原书图 45.1 将极半径极大值（顶点）误绘于正左方，与参数方程在 $\theta=0$ 时 $(x,y)=(l+h,0)$ 位于正右方的真实方位相差 $180^\circ$；本版依参数方程统一校准图 45.1～45.4 的坐标方位。
  \item \textbf{第 77 节 图 77.1（滑绳线平衡装置）}：原书图 77.1 在天花板线上标出曲线交点 $D$（正文并未提及该字母），并将定滑轮 $M$ 绘于 $D$ 点右侧，而正文推导则简记 $OM = l$。本版保留原图点位记号与定滑轮物理切线关系，并在随文脚注中说明此处“视滑轮中心与绳迹交点 $D$ 重合”的理想化近似。
  \item \textbf{第 95 节 图 95.1（对数螺线垂足曲线）}：原书图中将极角 $\theta$ 与 $\theta_1$ 绘作自极轴顺时针向下的小弧段，易在直观上引起 $|\theta_1|>|\theta|$ 的误解（实则顺时针角度均为负值，$\theta_1 = \theta - (\frac{\pi}{2}-\psi) < \theta < 0$）；本版适当调整螺线参数与角弧半径，使几何夹角与符号关系一目了然。
\end{enumerate}

\noindent\textbf{附记}：第 12 节图 12.2（微分三角形）中，原书将斜边 $ds$ 绘作微弧段；因弧长微分 $ds = \sqrt{dx^2+dy^2}$ 在几何上严格对应切线段（即微分直角三角形的斜边），本版依严谨数学概念改绘为直线斜边，全书其余古典数学符号规范均予忠实保留；所有物理光学与机械示意图（如图 4.1 探照灯遮蔽关系、图 5.1 卡塞格林望远镜实体镜片、图 6.1 杰尼西亚山洞、图 21.1 熊猫滚圆玩具、图 42.1 套藤圈等）亦均在重绘中依据物理与几何原理一并校准。
}
